\subsection{Wo die Ergebnisse der Sitzung in dieses Bild gehören}
Die Abschnitte~\ref{sec:physicalconstraints} und~\ref{sec:carriertraceclosure} erhalten die ursprünglichen Massen-, Flavor- und Alpha-Herleitungen samt ihrer heute ausgeschriebenen gemeinsamen Spurabbildung. Abschnitt~\ref{sec:operatorbridge} verbindet die Fuchs-Geometrie und den neutralen Nahtkandidaten mit tatsächlichen Quellenoperatoren. Die Neutrinopaare und ihre Norm-/Phasenabbildung stehen in Abschnitt~\ref{sec:spinliftsynthesis}; die Wandprüfung in Abschnitt~\ref{sec:wallaudit} korrigiert den behaupteten unbedingten Dirac-Zwang. Die Existenz eines geladenen Quellpaars ersetzt dabei nicht seine physische Masseneinsetzung.

Die Abschnitte~\ref{sec:masszeromode} und~\ref{sec:stateactionclosure} verbinden den vorhandenen schweren Neutrinoblock mit Masse, Rückwirkung und reduziertem Zustand. Die Clock-, Replica- und Gedächtnisresultate aus den Abschnitten~\ref{sec:nativeclosure}--\ref{sec:sharedtime} prüfen verschiedene Auslesungen derselben Quelle. Abschnitt~\ref{sec:inverseorigin} gewinnt aus dem markierten E8-Kandidaten Fermionenfeldregel, Vakuum und Konformalzeit gemeinsam zurück. Diese Ergebnisse sind Anschlüsse an die gemeinsame Aufgabe; keiner ersetzt allein die physische Gesamtrealisierung.

\subsection{Was als vollständige Lösung zählen würde}
Eine Lösung muss zugleich die vollständige Quelle aus den ursprünglichen Bedingungen auswählen, ihre gemeinsame Wirkungs- und Zustandsregel konstruieren und die vorhandenen physikalischen Auslesungen in einer konsistenten lokalen Raumzeit realisieren. Dabei müssen auch Kopplungen, Neutrinosektor, Gravitation und kosmologische Anfangs-/Transferbedingungen am selben Objekt hängen. Ein metrologischer Maßstab darf ausdrücklich kalibriert werden; zusätzliche dimensionslose Eingaben wie das bisher extern eingesetzte $\alpha_s(M_Z)$ müssen erkennbar bleiben.

Die zentrale Aufgabe ist deshalb eine \emph{gemeinsame konsistente Operator- und Prozessrealisierung mit begründetem Auslesewörterbuch}. Die vorhandenen Einzelgesetze und ihre gemischten Antworten müssen darin gemeinsam gelten. Ein erzeugendes Funktional ist eine mögliche Darstellung solcher Antworten; der originale Vertrag \file{FTRANSFER.GENERATING.01} legt es für vier bestimmte nachgelagerte Transfers fest. Abschnitt~\ref{sec:commonvariation} untersucht eine weitere bedingte Operatorprüfung. Keine dieser Formulierungen wählt die Herkunft bereits aus. Der globale Auswahlsatz in Abschnitt~\ref{sec:globaloriginselection} behandelt davor das ursprüngliche diskrete Gerüst mit einer ausdrücklich zusätzlichen Stromforderung.

\section{Der Ursprung: P1, P2 und die rekursive Rückbestätigung}

\subsection{P1: Orientierung, Rückkehr und positive Antworten}
P1 beginnt mit einem orientierten Seam und einem positiven beziehungsweise reflexionspositiven Kern. Die primitive Windung fixiert im verwendeten TFPT-Vertrag die topologische Normierung
\begin{equation}
c_3=\frac{1}{8\pi}.
\end{equation}
Eine positive Antwort bedeutet bildlich: Die aus der Quelle berechneten Normen dürfen nicht negativ werden. Orientierung bedeutet, dass Hin- und Rückweg sowie eine umgekehrte Windung nicht einfach dieselbe markierte Operation sind. ``Primitiv'' bezeichnet hier die ursprüngliche Windungs- beziehungsweise Unzerlegbarkeitsbedingung. Daraus folgt nicht ohne weiteren Satz die Perron--Frobenius-Primitivität irgendeines später eingeführten Vergröberungskanals.

\picturetext{Auf einem verdrehten Band kann man nach einer Runde am gleichen Ort sein und dennoch eine andere Orientierung mitbringen. Die Ortsanzeige allein beschreibt den Zustand nicht vollständig. Genau diese Art verlorener Information motiviert Double Cover, Phasenregister und Carry.}

Double Cover, Möbiusbild und Spinstruktur sind verwandte Motive, aber keine austauschbaren Beweise. Eine zweifache Überlagerung ist zunächst eine topologische Aussage; eine Spinordarstellung, eine Eichladung und eine tatsächlich messbare Interferenzphase benötigen zusätzlich die jeweils passende Wirkung auf Feldern und Zuständen.

\subsection{P2: Die markierte Aufteilung $3+2$}
P2 stellt die markierte Aufteilung des fünfdimensionalen Trägers bereit. Auf der fundamentalen Darstellung wird der Generator als
\begin{equation}
Y=\diag\left(-\frac13,-\frac13,-\frac13,\frac12,\frac12\right)
\end{equation}
geschrieben. Die Spurfreiheit ist exakt: $3(-1/3)+2(1/2)=0$. Bei der entsprechenden Trägernormierung entstehen $k_Y=5/3$ und die bekannte Hochskalenrelation $\sin^2\theta_W=3/8$. Das sind Aussagen an der markierten Trägerstruktur. Ihre Identifikation mit laufenden, renormierten Messgrößen erfordert den zugehörigen physikalischen Anschluss.

Die Sitzung hat das Vorzeichenproblem der neuen $X$-Felder ausdrücklich aufgelöst: Sie tragen $\overline{\mathbf5}\otimes\mathbf4$, also bezüglich P2 die entgegengesetzten fundamentalen Ladungen. Der Name ``Fünferblock'' darf dieses Dualisieren nicht verdecken.

\subsection{Anker, Familiengeometrie und Bootstrap}
Die ursprüngliche Ankerstruktur verwendet
\begin{equation}
a=(1,1,2),\qquad e_1=4,\quad e_2=5,\quad e_3=2,
\qquad p_n=2+2^n.
\end{equation}
Damit ergeben sich
\begin{equation}
p_1p_2p_3=4\cdot6\cdot10=240,\qquad p_4-p_3=18-10=8,
\qquad 240+8=248.
\end{equation}
Diese Zahlen passen zur Wurzelzahl, zum Rang und zur Dimension von $E_8$. Die tragfähige Identifikation kommt jedoch aus der tatsächlich konstruierten Gitter- und Gruppenwirkung, nicht allein aus diesen Gleichheiten.

Die vierfach punktierte projektive Gerade hat $H_1\cong\Z^3$; die ursprüngliche Familienlesart benutzt diese Rang-drei-Struktur. Ein Divisor vom Grad vier auf $\mathbb P^1$ besitzt fünf globale Schnitte. Die Markierungen $1,i,-1,-i$ geben das $\mu_4$-Register. Dies erklärt, weshalb Vierermarken, ein Fünferträger und drei unabhängige Schleifen in derselben Geometrie auftreten können, ohne dieselbe Art von Dimension zu bezeichnen.

Die Origin-Story und ihre Lean-Aussagen enthalten echte Rückbestätigungen: Unter den jeweiligen primitiven Determinanten-, Rang- und Yukawa-Voraussetzungen kehren $3+2$, die Ladungsrelationen und die entsprechenden Ränge wieder. Diese Rekursion darf nicht unterschlagen werden. Ihr logischer Inhalt ist aber eine Äquivalenz beziehungsweise ein Schluss innerhalb benannter Voraussetzungen. Ein Zyklus ``A liefert B, B liefert A'' beweist die Selbstkonsistenz von A und B; er beweist nicht von selbst, dass jede andere Ausgangsstruktur unmöglich ist.

\key{Die Bootstrap-Rückbestätigung ist ein struktureller Gewinn: Die Teile stützen dieselben Markierungen. Der nächste Ursprungssatz muss zeigen, welche konkurrierenden Quantisierungen oder Realisierungen durch genau diese Markierungen ausgeschlossen werden.}

\subsection{Was die erneute Prüfung des behaupteten Ursprungsabschlusses ergibt}\label{sec:originproofaudit}
Die stärkste archivierte Fassung minimiert auf zulässigen einseitigen spektralen Bordismen den Defekt
\begin{equation}
\mathfrak D(B)=\bigl(|\operatorname{SF}|,\rank_{\rm ess},\deg^+_{\det},h_\Sigma^{\rm red}\bigr).
\end{equation}
Die Wohlordnung dieser diskreten Werte liefert bei nichtleerer Klasse einen minimalen \emph{Wert}. Sie beweist nicht, dass genau ein vollständiger Operator diesen Wert annimmt. Das Original benennt die dafür notwendige Klassifikation der minimalen Faser selbst, führt sie im anschließenden Eindeutigkeitsbeweis aber nicht aus: \file{01_boundary_kernel_source.tex}, Zeilen 543--593 und 650--682, unter \file{_archive/tfpt-45/source_extracts/}. Die spätere Master-Barriere ist null für $\mathcal B\cong\mathcal B_{\min}$ und sonst unendlich. Sie setzt die ausgewählte Klasse damit schon ein. Eine rückgekoppelte Auswahl wäre grundsätzlich zulässig; dafür müssten unabhängig definierte Bedingungen tatsächlich gemeinsam gelöst und ihre Eindeutigkeit bewiesen werden.

Der konkrete ausführbare Abschluss \file{verification/v165_premise_a_closure.py} behauptet, aus dem DtN-/Calder\'on-Symbol $|k|$ folge bereits der freie Fermionen-Zweipunktkern. Der dafür ausgeführte Test prüft Rang-, Dimensions- und Indexgleichheiten. Sein Vorgänger \file{v162_seam_transport_identification.py} setzt die Rang-acht-Polarisation als $P=(I+iA_{16})/2$ ein und prüft ihre Projektoreigenschaft. Beide Rechnungen liefern keinen Intertwiner vom ursprünglichen Kragenoperator zu dieser Polarisation. Das Hauptsymbol allein bestimmt weder die vollständige Operatorwirkung und ihre Randbedingungen noch die behauptete Feldquantisierung. Der entsprechende \emph{Quellenabschluss} ist durch diese Tests daher nicht bewiesen.

Auch die originale Definition $Z_{\rm rel}=Z_{\rm adm}/Z_{\rm ref}[0]$ liefert keine zusätzliche Auswahl des Feldgesetzes. Sie normiert ein bereits gegebenes Funktional; die Reflexionspositivität seines CAR-Kerns ist im folgenden Lemma eine ausdrückliche Voraussetzung. Die Ableitungen von $\log Z_{\rm rel}$ definieren anschließend die Korrelatoren, aber noch nicht das zuvor benötigte $Z_{\rm adm}$ selbst. Maßgebliche Stellen sind \file{02_carrier_source.tex}, Zeilen 2257--2311, und \file{04_qft_source.tex}, Zeilen 1499--1515.

Der noch mögliche geometrische Reparaturansatz ist ebenfalls präzise begrenzt. Die skalare Flachheitsbedingung für $|D_\theta|+M_f$ wählt keine Verbindung auf dem gesamten $V_{16}$ aus. Der vorhandene Master variiert $(\alpha,\chi_{\rm geo},\delta_{\rm ph},\rho_{\rm vac})$, nicht die blockmischende Verbindung dieses Trägers. Eine flache Gesamtverbindung wäre geometrisch möglich; ihre relative Lage zu den markierten Teilbündeln müsste zusätzlich aus den ursprünglichen Bedingungen folgen. Sie hier einzusetzen wäre eine neue Annahme.

\textbf{Reichweite dieses Befunds:} Der behauptete Quellenabschluss ist an einer bestimmten Beweisstelle nicht getragen. Es ist damit weder die Unmöglichkeit von TFPT noch die Existenz zweier vollständiger Modelle unter sämtlichen gemeinsamen Bedingungen bewiesen. Der fehlende Satz bleibt die Auswahl des vollständigen markierten lokalen Operators und seines Quellfunktionals aus den Ursprungsbedingungen. Weitere Rechnungen innerhalb der schon gewählten E8-Quelle schließen diesen Satz nicht nachträglich.

\subsection{Drei verschiedene Eindeutigkeitsfragen}
Die Rückbestätigung des diskreten Gerüsts, der stationäre Zustand eines gegebenen Transfers und die Auswahl eines vollständigen physischen Gesetzes haben verschiedene Beweisgegenstände. Die erneute Originalprüfung macht die Verwechslung genau lokalisierbar: \file{v56_unique_attractor.py} setzt eine Matrix
\[
T=V\,\diag\bigl(1,(2/3)^6,(1/3)^6\bigr)V^{-1}
\]
ein und prüft ihre führende Richtung sowie Konvergenz. Der erfolgreiche Test beweist keine Eindeutigkeit der Auswahl von $T$ aus allen ursprünglichen Bedingungen.

\file{v353_selfloop_capstone.py} prüft mehrere diskrete Rückrelationen; für die weitergehende Aussage ``zero free dials'' wird im Code jedoch \file{free_dials=0} gesetzt. Zwei abschließende weitreichende Aussagen werden mit \file{check(..., True)} registriert. Die frischen Läufe bestanden fünf beziehungsweise acht interne Prüfungen (v353 beziehungsweise v56). Das erhält die tatsächlich geprüften Rückrelationen, macht aus den weitergehenden gesetzten Aussagen aber keinen Beweis.

Damit ist weder die vollständige TFPT widerlegt noch ihre Mehrdeutigkeit unter sämtlichen Bedingungen nachgewiesen. Insbesondere respektiert ein beliebig konjugierter Operator mit denselben Eigenwerten nicht automatisch die bereits festgelegten Feldprodukte, Phasen und gemischten Antworten. Der noch benötigte Auswahlbeweis betrifft genau diese vollständige markierte Struktur.

\section{Vom Double Cover zum $E_8$-Träger und zum Standardmodell}

\subsection{Die Verklebung ist eine konkrete Ladungsrechnung}
Die alte Architektur benutzt die $D_5$-Spinorstruktur
\begin{equation}
S^+\cong\Lambda^{\mathrm{even}}\C^5,\qquad \dim S^+=16,
\end{equation}
und den $A_3$-Registerpartner. Die Diskriminantengruppen der verwendeten Gitter tragen passende $\Z_4$-Klassen. Die Normbeiträge $5/4$ und $3/4$ addieren sich zu $2$. Die Index-vier-Verklebung liefert das gerade unimodulare Rang-acht-Gitter $E_8$; seine vier ursprünglichen Wurzelklassen haben Größen $52+64+60+64$.

Die zugehörige Zerlegung der adjungierten Darstellung lautet
\begin{equation}
\mathbf{248}=(\mathbf{45},\mathbf1)\oplus(\mathbf1,\mathbf{15})
\oplus(\mathbf{16},\mathbf4)\oplus(\overline{\mathbf{16}},\overline{\mathbf4})
\oplus(\mathbf{10},\mathbf6).
\end{equation}
Die Gleichung verbindet eine tatsächliche Gruppenwirkung mit Ladungssektoren. Die oft verführerische Gleichheit $16=16$ wäre allein unzureichend: Feldkomponenten eines geometrischen Operators sind nicht automatisch die sechzehn chiralen Materiezustände einer Generation.

\subsection{Was die Architektur liefert und was die Physik zusätzlich verlangt}
Die Architektur liefert einen gemeinsamen Organisationsraum für interne Ladungen, Register und mögliche Materieblöcke. Anomaliefreiheit ist eine notwendige Konsistenzbedingung. Sie wählt nicht automatisch eine Wechselwirkung, entkoppelt keine Spiegelmaterie und präpariert keinen physischen Zustand.

Für den Standardmodellanschluss müssen daher die tatsächlichen Ladungsdarstellungen, die chiral wirksamen Freiheitsgrade, die lokalen Wechselwirkungen, der Zustand und die Raumzeitdynamik zusammenkommen. Die Sitzung benutzt die $E_8$-Kaskade ausdrücklich nicht als vorausgesetzten universellen Entstehungsmechanismus. Sie bleibt eine mögliche Geschichte unter zusätzlichen Voraussetzungen; der direkte Gitterquellenweg benötigt sie nicht.

\section{Die verzahnten Clocks, Flavor und $\alpha$ im Gesamtgefüge}

\subsection{Mehrere Uhren können dieselbe Struktur ablesen}
Die ursprünglichen Clock-Strukturen enthalten
\begin{equation}
h(E_8)=30=2\cdot3\cdot5,\qquad \varphi(30)=8,
\end{equation}
mit den Exponenten $1,7,11,13,17,19,23,29$. Daneben steht die vierfache Markierung. $\mu_4$ ist keine Untergruppe eines zyklischen 30-Taktes. Der tatsächliche Anschluss verwendet unter anderem eine Wirkung der Form
\begin{equation}
G C_5G^{-1}=C_5^2,\qquad G^4=I.
\end{equation}
Die komplexe Reflexionsgruppe $G_{31}$ hat Ordnung $46080$, sechzig Reflexionen und Invariantengrade $8,12,20,24$. Diese Angaben erklären gemeinsame diskrete Symmetrien; sie stellen noch keine Uhr mit einer physisch festgelegten Sekundenskala dar.

In einer verwendeten Clock-Zerlegung gilt $c=J\sigma_F$ von Ordnung zwölf, mit $\sigma_F=c^4$ und $J=c^9$. Zwei Operatoren gleicher Ordnung zwölf müssen nicht konjugiert sein und nicht gleich auf den geladenen Feldern wirken. In der Sitzung wurde deshalb die tatsächliche Wirkung statt nur der Periodenlänge verglichen.

\subsection{Die bereits vorhandene physikalische Herleitungskette}\label{sec:physicalconstraints}

\key{Die Aussage ``Massen, Flavor, Yukawa und $\alpha$ müssen noch hergeleitet werden'' war zu pauschal. Die Originale enthalten konkrete Formeln, Operatoren und Begründungen. Die präzise Frage lautet: Realisiert ein und derselbe ursprüngliche Quellprozess diese vorhandenen Gesetze mit denselben Markierungen, Normierungen und Zuständen?}

Die erneute Prüfung benutzt die aktuellen PDF-Fassungen und ihre TeX-Quellen, die Originalverifier und die späteren Quellverträge. Im Standardmodell-PDF stehen die geladenen Leptonen auf Seite 8, die Quarkverhältnisse auf Seiten 10--11 und der CKM-Anschluss auf Seite 29. Die unterschiedlichen Beweisstärken sind in den Originalen selbst markiert; sie werden hier weder pauschal abgesenkt noch zusammen zu einem bereits bewiesenen physikalischen Abschluss erklärt.

\begin{longtable}{L{0.19\linewidth}L{0.45\linewidth}L{0.27\linewidth}}
\toprule
Bereich & Vorhandenes Ergebnis & Genaue Reichweite\\\midrule\endhead
Geladene Leptonen & Potenzen von $\phi_0$ und Koeffizienten $16/7,4/3,7/6$ aus der Familienmonodromie & Exakte Quellenverhältnisse; Legzuordnung, Maßstab und RG-Abbildung mitführen\\
Quarks & Exakte Verhältnisse aus Plücker-/Compiler-Auslesungen; Starrheit auf dem diskreten Selektorstratum & Kein weiterhin offenes vollständiges Hitchin-Problem für diese Verhältnisse; absolute Amplituden gesondert prüfen\\
Yukawa und CKM & Gemeinsame $K/L$-Hierarchie, konkrete Mischungswinkel und führende CP-Phase & Höhere CP-Korrektur und vollständiger gemeinsamer Operatoranschluss getrennt\\
Neutrinos & Endlicher Dirac-/Majoranaoperator, Seesaw, PMNS-Kanäle, Scalaron-Skalenkandidat & Nicht alle Zweige sind derselbe Ansatz; schwerer Sektor und Zerfallsquelle müssen gemeinsam passen\\
Feinstruktur\-konstante & Konkrete Determinantenantwort, Potential und eindeutige positive Nullstelle & Exakte Lösung der angegebenen Gleichung; Herkunft als volle physische Determinantenvariation gesonderter Vertrag\\
Gravitation und Skalen & Stationäres Planckpotential, EW-/Vakuumenergie-Elimination und Scalaronroute & Bedingte gemeinsame Funktional- und Matchingannahmen; keine beliebig unabhängigen Skalen\\
\bottomrule
\end{longtable}

\subsection{Massen und Yukawa-Kopplungen aus derselben markierten Struktur}
Mit Zeilen $u,d,e$ lauten die ursprünglichen Operatoren beziehungsweise markierten Auslesetabellen
\begin{equation}
R=\begin{pmatrix}1&3&0\\1&5&2\\2&5&3\end{pmatrix},\quad
K=\begin{pmatrix}4&2&0\\4&3&2\\5&3&2\end{pmatrix},\quad
L=\begin{pmatrix}7&3&0\\7&5&2\\8&5&3\end{pmatrix}.
\end{equation}
Es gilt $K=R+Q\Sigma$, $L=R+Q(I+\Sigma)$ und $\sum_{f,j}L_{f,j}=40$. Der Seed und das gemeinsame Yukawa-Gesetz sind
\begin{align}
c_3&=\frac1{8\pi},\qquad \phi_0=\frac1{6\pi}+48c_3^4,
\qquad \lambda_Y=\sqrt{\phi_0(1-\phi_0)},\\
\hat y_{f,j}&=\lambda_Y^{L_{f,j}}\Lambda_{f,j}
=\pi c_{f,j}\phi_0^{K_{f,j}},\qquad
\hat m_{f,j}=\frac{v_{\rm geo}}{\sqrt2}\hat y_{f,j}.
\end{align}
Hüte kennzeichnen die Quellenwerte. Die Standard-RG-/Schwellenabbildung zu laufenden oder Polmassen ist ein eigener, benannter Rechenschritt. Sie macht die Quellenherleitung nicht wertlos; Quellenmassen dürfen aber nicht ungeprüft mit Messmassen gleichgesetzt werden.

Besonders konkret ist die Leptonenrechnung. Die ursprüngliche Familienmonodromie $M_F$ aus v117/v118 erfüllt
\begin{equation}
\det(I-tM_F)=1-t^3,\qquad d_-:=\det(I-M_F/2)=7/8,
\quad d_+:=\det(I+M_F/2)=9/8.
\end{equation}
Damit ergeben sich
\begin{equation}
c_e=\frac4{2d_-}=\frac{16}7,\quad
c_\mu=\frac3{2d_+}=\frac43,\quad
c_\tau=\frac{4d_-}3=\frac76.
\end{equation}
Die Vereinigung der Spektren von $M_F$ und $-M_F$ liefert die sechs Phasen: Familienrotation und Blattvorzeichen gehören zusammen. Die Monodromie wird hier $M_F$ genannt, damit sie nicht mit der anderswo ebenfalls $M_0$ genannten Gramform verwechselt wird. Der Originaltest setzt die etablierten Familienlegs ein; er leitet ihre physische Zuordnung nicht nochmals aus der Determinante her.

Folglich gilt exakt
\begin{equation}
(\hat m_e,\hat m_\mu,\hat m_\tau)
=\frac{v_{\rm geo}\pi}{\sqrt2}
\left(\frac{16}7\phi_0^5,\frac43\phi_0^3,\frac76\phi_0^2\right).
\end{equation}
Dies liefert zwei besonders klare gemeinsame Bedingungen:
\begin{equation}\label{eq:physicalleptonelimination}
\boxed{\frac78\frac{\hat m_\mu}{\hat m_\tau}=\phi_0,
\qquad
\frac{\hat m_e\hat m_\tau^2}{\hat m_\mu^3}=\frac{21}{16}.}
\end{equation}
Die zweite Gleichung eliminiert sowohl den gemeinsamen Maßstab als auch den Seed. Sie ist eine direkte algebraische Konsequenz der vorhandenen Formeln, kein neuer unabhängiger Messbefund und keine beanspruchte mathematische Neuheit.

Auch die Quarkverhältnisse sind konkret:
\begin{equation}
\frac{\hat m_u}{\hat m_d}=\frac{55}{117},\quad
\frac{\hat m_c}{\hat m_s}=\frac{34}{47}\phi_0^{-1},\quad
\frac{\hat m_t}{\hat m_b}=\frac3{26}\phi_0^{-2}.
\end{equation}
Die Originalkennungen \file{FLAV.RIGID.01/02} erhalten ihre Bedeutung: Diese Verhältnisse sind auf dem hergeleiteten diskreten Selektorstratum starr. Die absolute Quarknormierung folgt nicht schon daraus, dass eine Prüfroutine eingesetzte Amplitudenprodukte wieder zurückrechnet. Genau diese stärkere Behauptung muss von den tatsächlich geschlossenen Verhältnissen getrennt bleiben.

Für CKM stehen im Original
\begin{equation}
s_{12}=\lambda_Y,\qquad s_{23}=\frac{\phi_0}{1+\lambda_Y},\qquad
s_{13}=\frac{\lambda_Y^3}{3},\qquad
\delta_{\rm CKM}=\frac\pi3+3\lambda_Y^2+\Delta_\triangle.
\end{equation}
Die führende Phase ist berechnet; $\Delta_\triangle$ bezeichnet die ausdrücklich noch nicht geschlossene höhere Korrektur. Der Koeffizient drei darf nicht nach einem günstigeren Datenvergleich ausgetauscht werden.

\subsection{Warum $48$, $40$ und $41$ eine gemeinsame Einschränkung sind}
Die Kopplung zwischen den beiden Herleitungswegen ist genauer als eine Zahlenähnlichkeit. Auf einer ursprünglichen Familie gilt
\begin{equation}
\tr_{16}Y^2=\frac{10}3,\qquad
\Omega_{\rm adm}=3\cdot16=48,\qquad
\gamma=\tr_{5}Y^2=\frac56.
\end{equation}
Ungewichtet ergeben die drei Familien den Punkturnachtrag $48c_3^4$. Mit Fermion-/Skalargewicht und Eichnormierung folgt dagegen
\begin{equation}
b_1=\frac35\left(\frac23\,3\,\frac{10}3+\frac13\,\frac12\right)
=\frac{41}{10},\qquad
\sum L=48\gamma=40,\quad \sum L+N_\Phi=41.
\end{equation}
Der Higgs trägt den letzten Beitrag bei. Ein ungeladenes rechtshändiges Neutrino trägt nicht zur $Y^2$-Spur bei, gehört aber zum ungewichteten Zustandsbestand. Es darf daher nicht aus der 48-Zählung entfernt werden, nur weil die elektromagnetische Spur seine Entfernung nicht bemerkt. $5/3$ ist der Eich-Einbettungsindex; $5/4=(3/2)(5/6)$ ist die Nahtsuszeptibilität. Beide haben unterschiedliche Rollen.

Der elektromagnetische Ansatz enthält bereits einen tatsächlichen Operator:
\begin{align}
U_\Sigma^{\min}&=|\psi_\Sigma\rangle\langle\psi_\Sigma|,
\quad \norm{\psi_\Sigma}=1,\qquad q=48c_3^4e^{-2\alpha},\\
\mathcal W_\Sigma&=-\frac54\tr\log(I-qU_\Sigma^{\min})
=-\frac54\log(1-q),\\
\phi_s(\alpha)&=\frac1{6\pi}+q(1-q)^{-5/4}.
\end{align}
Rang eins bezeichnet den ausgewählten primitiven Transferkanal, nicht die gesamte Quelle. Der unresummierte Flavorseed $\phi_0$ und $\phi_s(\alpha)$ sind verwandte, aber unterschiedliche Antworten.

Auch das Potential ist ausdrücklich vorhanden:
\begin{align}
\Gamma_{U(1)}(\alpha)&=\frac{\alpha^4}{4}-\frac23c_3^3\alpha^3
-8b_1c_3^6\int_0^\alpha\log\bigl(1/\phi_s(a)\bigr)\,da,\\
\partial_\alpha\Gamma_{U(1)}&=
\alpha^3-2c_3^3\alpha^2-\frac45\,41c_3^6\log(1/\phi_s)=0.
\end{align}
Die positive Lösung ist $\alpha^{-1}=137.0359992168407\ldots$. Die neue Prüfung bestätigt diese interne Schließung. Der Originalvertrag \file{ALPHA.QUILLEN.EXACT.01} fragt genauer nach ihrer Herkunft als vollständige relative Determinantenvariation mit dem richtigen lokalen Gegenbeitrag. Ein vorhandenes Potential zu verneinen wäre falsch; es allein durch Integration der Zielgleichung physikalisch herzuleiten ebenso.

\picturetext{Einmal zählt man alle Plätze in einem Fahrzeug. Ein zweites Mal zählt man dieselben Plätze mit Gewichten nach ihrer Ladung. Die beiden Summen sind verschieden, aber man kann die Sitzordnung nicht in jeder Rechnung neu erfinden. Genau so müssen $48$, $41$, die Flavorzuordnung und die Nahtantwort aus demselben Bestand kommen.}

\subsection{Der Flavor-Determinant und die Eichspur aus gemeinsamen Ladungsdaten}\label{sec:carriertraceclosure}
\key{Der ursprüngliche Flavorpol ist bereits aus einer gemeinsamen Wirkung hergeleitet. Die heutige Verbindung erklärt ihre genaue Einsetzung mit denselben $S^+$-Ladungsdaten und dem vorhandenen gedoppelten Familienweg. Entscheidend sind die Art der Projektion und die Art der Spur.}

Das Original definiert in \file{03_em_flavor_source.tex:823-997} unter
\file{_archive/tfpt-45/source_extracts/}
\begin{align}
D_{y,\delta}&=yI_6-\delta U_6,\\
\Gamma_{\rm tr}(\delta;\varepsilon)
&=-\sum_{y\in\{1,2/3,1/3\}}
\log\det(D_{y,\delta}^{\dagger}D_{y,\delta}+\varepsilon^2I_6).
\end{align}
Die drei $y$ sind die Hyperladungsbeträge von $e^c,u^c,d^c$ auf
$S^+=\Lambda^{\rm even}(\C^3\oplus\C^2)$, siehe dort Zeilen 1779--1800.
Diese Felder sind \emph{schwache} Singuletts. Die beiden Quarktypen sind keine Farbsinguletts.
Das sechsdimensionale Transportregister enthält bereits die drei Familienkomponenten und die zwei Faserzugänge des ursprünglichen Weges.

Auf dem einen Familienpaket definiert die Auswahl der geladenen schwachen Singuletts
\[
V_{\rm cusp}=e^c\oplus u^c\oplus d^c,\qquad
(\rank P_e,\rank P_u,\rank P_d)=(1,3,3).
\]
Der konkrete gemeinsame interne Operator und die normalisierte Farbspur lauten
\begin{align}
D_{\rm cusp}&=|Y|_{V_{\rm cusp}}\otimes I_6-\delta I_7\otimes U_6,\\
W_{\rm cusp}&=P_e+\tfrac13P_u+\tfrac13P_d,\\
\Gamma_{\rm tr}(\delta;\varepsilon)
&=-\Tr_{V_{\rm cusp}\otimes\C^6}
\left[(W_{\rm cusp}\otimes I_6)
\log(D_{\rm cusp}^{\dagger}D_{\rm cusp}+\varepsilon^2I)\right].
\end{align}
Die Identität folgt aus dem Funktionalkalkül auf den drei Hyperladungsblöcken:
Jeder Block ist $I_{d_y}\otimes(D_{y,\delta}^{\dagger}D_{y,\delta}+\varepsilon^2I_6)$,
und $\Tr(P_y/d_y)=1$. Es werden alle Farbkomponenten behalten und ihre Antworten je Multiplet gemittelt.
Das ist eine exakte Umschreibung der ursprünglichen strukturellen Wirkung.
Die gewöhnliche ungewichtete Spur auf allen $42=7\cdot6$ Komponenten würde stattdessen die Gewichte $(1,3,3)$ erzeugen und ist ein anderes Funktional.
Der volle Materieträger $S^+\otimes F_3$ besitzt weiterhin $48$ Komponenten; hier wird keine zusätzliche Familie hinzutensoriert. Der $42$-dimensionale Transportraum ist kein mit ihm identischer Hilbertraum: Er verwendet die geladenen schwachen Singuletts und den bereits vorhandenen Familien-/Zweifaserfaktor. Die vollständige physische Abbildung zwischen diesen Registern gehört zum noch zu beweisenden Nahtanschluss.

\textbf{Die Quelle erhält die Quarks.} Die Originaldefinition des lokalen Netzes benutzt die Einschränkung invariant erhaltender Operatoren auf $\mathcal H_{\rm adm}\subset\mathcal H_{\rm OS}$; sie warnt ausdrücklich vor einer Definition durch globale Kompression
(\file{01_boundary_kernel_source.tex:315-342}). Die hadronische Algebra besteht aus center-neutralen \emph{Wörtern}
(\file{04_qft_source.tex:877-908}). Quarkpaare und geschlossene Quarkantworten bleiben daher verfügbar.
Die entsprechende Unterscheidung auf dem endlichen Träger ist
\[
P_c=\tfrac13(I+Z_c+Z_c^2),\qquad
\mathcal E_c(A)=\tfrac13\sum_{j=0}^2 Z_c^j A Z_c^{-j}.
\]
$P_c$ projiziert Vektoren, $\mathcal E_c$ bildet center-neutrale Operatoren.
Wegen $[Y^2,Z_c]=0$ gilt $\mathcal E_c(Y^2)=Y^2$; die volle Spur $10/3$ je Familie und $b_1=41/10$ bleiben erhalten.
Die \emph{unzulässige andere Operation}, zuerst alle farbigen Einteilchenfelder wegzulassen, ergäbe nur $12$ statt $48$ Komponenten und $b_1=19/10$. Dieser Gegenfall ist keine Behauptung über die im Original tatsächlich definierte Wortalgebra.
Center-Neutralität allein ist wiederum nicht vollständige lokale $SU(3)$-Eichinvarianz; die hier benutzten Hyperladungsprojektoren sind zusätzlich unter der ganzen internen Standardmodellwirkung invariant.

\textbf{Der Umlauf bleibt bei der Reduktion vollständig erhalten.}
Für die originale Matrix $U_6=\diag(M_F,-M_F)$ liefert
\[
v=\frac1{\sqrt2}(1,0,0,1,0,0)^{\mathsf T},\qquad
K=(v,U_6v,\ldots,U_6^5v)
\]
den expliziten unitären Basiswechsel $K^\dagger U_6K=C_6$, wobei $C_6$ die zyklische Verschiebung ist.
Das wurde an derselben ursprünglichen Monodromiematrix exakt geprüft; es ist kein Ersatz nur durch passende Eigenwerte.
Behalte in jedem Zyklus einen Platz und eliminiere die übrigen fünf mit einem Schurkomplement. Dann gilt
\begin{equation}\label{eq:flavorschur}
\det E_y=y^5,\qquad
S_y=y-\frac{\delta^6}{y^5},\qquad
(D_{y,\delta}^{-1})_{00}=S_y^{-1},\qquad
\det D_{y,\delta}=y^5S_y=y^6-\delta^6.
\end{equation}
Die fünf verborgenen Plätze bleiben als vollständiger Umlauf $\delta^6/y^5$ wirksam.
Die naive Ein-Platz-Kompression von $D^\dagger D$ wäre dagegen nur $y^2+\delta^2$ und würde diese Antwort verlieren.

Für $\varepsilon=0$ folgt der bereits originale vollständige Abschluss
\begin{align}
P(z)&=(z-1)(z-64/729)(z-1/729),\\
\Gamma_{\rm tr}(\delta;0)&=-2\log|P(\delta^6)|,\\
\Gamma_{\rm tr}'(\delta)&=-12\delta^5\frac{P'(\delta^6)}{P(\delta^6)},\\
\delta_*&=\left(\frac{794-7\sqrt{9961}}{2187}\right)^{1/6}
=0.593277280133224001\ldots .
\end{align}
In der ursprünglichen Kammer $1/3<\delta<2/3$ ist dies die einzige stationäre Stelle;
$\Gamma_{\rm tr}''(\delta_*)=417.715131947579643\ldots>0$.
Die Kammer ist wesentlich, da die Kubik im anderen Wurzelzwischenraum ebenfalls eine stationäre Stelle liefert.
Für $\varepsilon>0$ bleibt die vollständige regulierte Determinante maßgeblich.
Die Betragsquadrate bestimmen zudem nicht allein die Phase einer gemeinsamen Quillen-/Eta-Linie.

\picturetext{Man lässt fünf Stationen eines Rundwegs aus der Zeichnung verschwinden. Ihr Einfluss verschwindet nicht: Er erscheint als Rückweg am sechsten Platz. Diese richtige Vergröberung erklärt zugleich, warum die kleine Kubik die stationäre Antwort des ganzen Transports behält.}

\textbf{Gesamtbedeutung und Grenze.} Die Träger-, Flavor- und Eichformeln lassen sich nun ohne neue Felder und ohne widersprüchliche Farbprojektion an dieselbe markierte Darstellung anschließen.
Die normalisierte Multipletspur ist dabei ausdrücklich sichtbar. Ihre Auswahl aus \emph{derselben physischen Wirkung}, welche die volle Eichspur erzeugt, bleibt nachzuweisen; eine verfügbar gemachte Spur ist noch kein Beweis ihrer dynamischen Auswahl.
Ebenso fehlt der vollständige Intertwiner zum ursprünglichen Nahtoperator mit lokaler Raumzeit.
Die Rechnung ist in \file{tfpt_explorer/source_charge_response.py} implementiert. Acht gezielte Tests bestanden: darunter die direkte regulierte $42$-dimensionale Matrixlogarithmus-Rechnung, der echte Zyklus-Intertwiner, die erhaltene Resolventenantwort und die falsche Einteilchenprojektion als Gegenkontrolle.

\subsection{Die gemeinsame Variationsidentität statt einzelner Extremwerte}\label{sec:commonvariation}
Eine bisher verfolgte Prüfroute untersucht die ursprüngliche gemeinsame Operatorfamilie $D_\xi$ mit
$\xi=(\alpha,\chi_{\rm geo},\delta_{\rm ph},\rho_{\rm vac})$.
Die Masterformel in \file{02_carrier_source.tex:3021-3145} muss ihre verschiedenen Reduktionen aus \emph{dieser einen Familie} erzeugen.
Der heute bewiesene Trägeranschluss liefert die richtigen Räume und Spurgewichte; die bereits berechneten neutralen, geladenen und gemischten E8-Antworten liefern die entsprechenden Vergleichsgrößen.

Für eine glatte invertierbare endliche Operatorfamilie gilt beispielsweise bei $\Gamma_{\det}=-\log\det D$
\begin{equation}\label{eq:mixedvariation}
\boxed{\partial_\alpha\partial_\delta\Gamma_{\det}
=\Tr(D^{-1}D_\alpha D^{-1}D_\delta)
-\Tr(D^{-1}D_{\alpha\delta}).}
\end{equation}
Dies folgt unmittelbar durch Ableiten von $D^{-1}$; keine neue Dynamik wird eingesetzt.
In der relativen, zeta-regulierten oder Pfaffian-Wirkung müssen zusätzlich Referenz-, Kontakt-, Phasen- und gegebenenfalls Halbierungsfaktoren korrekt mitgeführt werden.
Die behauptete sektorale Separation des Masters verlangt für die \emph{gesamte} Wirkung eine entsprechende Nullidentität der gemischten Ableitung.
Kommutierende Selektoren allein beweisen diese Identität nicht.
Umgekehrt widerlegt ein nichtverschwindender gemischter Quellenkorrelator sie auch nicht: Der Kontaktterm und die übrigen Wirkungsbeiträge können sie kompensieren.

Die ursprüngliche Flavorseite hat $\partial_\delta D_{y,\delta}=-U_6$; die Alpha-Seite besitzt ihren Rang-eins-Port und $q(\alpha)=48c_3^4e^{-2\alpha}$.
Gesucht ist ihre \emph{gemeinsame} Einsetzung einschließlich der Normalisierung $W_{\rm cusp}$, der gemischten Vertices und der geometrischen Variation.
Das ist eine konkrete gemeinsame Herleitungsaufgabe. Eine bloße Summe nachträglich integrierter skalarer Zielgleichungen liefert diesen Operatoranschluss nicht.
Die Verträge sind getrennt zu lesen: \file{ALPHA.QUILLEN.EXACT.01} und \file{GAUGE.DETLINE.FIXPOINT.01} betreffen bestimmte elektromagnetische beziehungsweise Eich-Determinantenanschlüsse. \file{FTRANSFER.GENERATING.01} verlangt demgegenüber, unter der vorausgesetzten vierdimensionalen Theorie mit kontrolliertem Kontinuum, einen gemeinsamen Erzeuger für die vier nachgelagerten Transfers. Die Gleichsetzung dieses Erzeugers mit der hier betrachteten Operatorfamilie folgt aus keiner dieser Benennungen. Die Variationsidentität ist ein Werkzeug für einen belegten Anschluss; sie wählt den fundamentalen Ansatz nicht aus.

\subsection{Die bekannten Formeln schränken den Ursprung rückwärts ein}
Auf dem verwendeten kleinen Seedzweig folgt bereits
\begin{equation}
\phi_0=\frac78\frac{\hat m_\mu}{\hat m_\tau}
=\frac{1-\sqrt{1-4s_{12,\rm CKM}^2}}2.
\end{equation}
Die ursprünglichen weiteren, jeweils bedingt physisch gelesenen Kanäle lauten
\begin{equation}
\phi_0=e^{5/6}\sin^2\theta_{13,\rm PMNS}
=4\pi\beta_{\rm rad}
=\frac{4\pi}{4\pi-1}\Omega_b.
\end{equation}
Diese Gleichheiten verknüpfen dieselben Quellenformeln; sie sind kein heute neu ausgeführter globaler Fit an Beobachtungsdaten. Ihre Abhängigkeiten müssen mitgeführt werden, statt jede Auslesung als unabhängige Bestätigung zu zählen.

Die Rückrichtung ist dennoch streng: Für $c>0$ ist
\begin{equation}
\phi(c)=\frac43c+48c^4,\qquad \phi'(c)=\frac43+192c^3>0.
\end{equation}
Ein konsistenter positiver Flavorseed wählt somit innerhalb dieser bereits festgelegten Funktionsklasse genau ein $c$. Zusammen mit dem Budget $41$ ist dann auch die angegebene $\alpha$-Schließung festgelegt. Das eliminiert unabhängige Nachjustierungen zwischen den Sektoren. Es beweist keine Eindeutigkeit jedes vollständigen Quelloperators, der diese endlich vielen Antworten reproduziert.

